%%% Hlavní soubor. Zde se definují základní parametry a odkazuje se na ostatní části. %%%

%% Verze pro jednostranný tisk:
% Okraje: levý 40mm, pravý 25mm, horní a dolní 25mm
% (ale pozor, LaTeX si sám přidává 1in)
\documentclass[12pt,a4paper]{report}
\setlength\textwidth{145mm}
\setlength\textheight{247mm}
\setlength\oddsidemargin{15mm}
\setlength\evensidemargin{15mm}
\setlength\topmargin{0mm}
\setlength\headsep{0mm}
\setlength\headheight{0mm}
% Doporučená úprava zmiňuje řádkování 1.5, to ale není pro TeX relevantní.
% \openright zařídí, aby následující text začínal na pravé straně knihy
\let\openright=\clearpage

%% Pokud tiskneme oboustranně:
% \documentclass[12pt,a4paper,twoside,openright]{report}
% \setlength\textwidth{145mm}
% \setlength\textheight{247mm}
% \setlength\oddsidemargin{14.2mm}
% \setlength\evensidemargin{0mm}
% \setlength\topmargin{0mm}
% \setlength\headsep{0mm}
% \setlength\headheight{0mm}
% \let\openright=\cleardoublepage

% Přepneme na českou sazbu
\usepackage[czech]{babel}
\usepackage[IL2]{fontenc}

%% Použité kódování znaků: obvykle latin2, cp1250 nebo utf8:
\usepackage[utf8]{inputenc}

%% Ostatní balíčky
\usepackage{graphicx}
\usepackage{amsthm}

%%% Údaje o práci

% Název práce v jazyce práce (přesně podle zadání)
\def\NazevPrace{Název práce}

% Název práce v angličtině
\def\NazevPraceEN{Name of thesis}

% Jméno autora
\def\AutorPrace{Jméno Příjmení}

% Rok odevzdání
\def\RokOdevzdani{ROK}

% Název katedry nebo ústavu, kde byla práce oficiálně zadána
% (dle Organizační struktury MFF UK, případně plný název pracoviště mimo MFF)
\def\Katedra{Název katedry nebo ústavu}
\def\KatedraEN{Name of the department}

% Jedná se o katedru nebo o ústav?
\def\TypPracoviste{Katedra}

% Vedoucí práce: Jméno a příjmení s~tituly
\def\Vedouci{Vedoucí práce}

% Pracoviště vedoucího (opět dle Organizační struktury MFF)
\def\KatedraVedouciho{katedra}
\def\KatedraVedoucihoEN{department}

% Studijní program a obor
\def\StudijniProgram{studijní program}
\def\StudijniObor{studijní obor}

% Nepovinné poděkování (vedoucímu práce, konzultantovi, tomu, kdo
% zapůjčil software, literaturu apod.)
\def\Podekovani{
Poděkování.
}

% Abstrakt (v rozsahu cca 80-200 slov; nejedná se o zadání práce)
\def\Abstrakt{
Abstrakt.
}
\def\AbstraktEN{
Abstract.
}

% 3 až 5 klíčových slov, každé uzavřeno ve složených závorkách
\def\KlicovaSlova{%
{klíčová} {slova}
}
\def\KlicovaSlovaEN{%
{key} {words}
}

%% Balíček hyperref, kterým jdou vyrábět klikací odkazy v PDF,
%% ale hlavně ho používáme k uložení metadat do PDF (včetně obsahu).
\usepackage[pdftex,unicode]{hyperref}   % Musí být za všemi ostatními balíčky
\hypersetup{breaklinks=true}
\hypersetup{pdftitle={\NazevPrace}}
\hypersetup{pdfauthor={\AutorPrace}}
\hypersetup{pdfkeywords=\KlicovaSlova}

%% Titulní strana a různé povinné informační strany
\input titulka.tex

%%% Strana s automaticky generovaným obsahem bakalářské práce. U matematických
%%% prací je přípustné, aby seznam tabulek a zkratek, existují-li, byl umístěn
%%% na začátku práce, místo na jejím konci.

\openright
\pagestyle{plain}
\setcounter{page}{1}
\tableofcontents

%%% Jednotlivé kapitoly práce jsou pro přehlednost uloženy v samostatných souborech
\chapter*{Úvod}
\addcontentsline{toc}{chapter}{Úvod}

Následuje několik ukázkových kapitol, které doporučují, jak by se
měla bakalářská práce sázet. Primárně popisují použití \TeX{}ové
šablony, ale obecné rady poslouží dobře i~uživatelům jiných
systémů.

\chapter{N�zev prvn� kapitoly}

\section{N�zev prvn� podkapitoly v prvn� kapitole}

\section{N�zev druh� podkapitoly v prvn� kapitole}


\chapter{Název druhé kapitoly}

\section{Název první podkapitoly v druhé kapitole}

\section{Název druhé podkapitoly v druhé kapitole}



% Ukázka použití některých konstrukcí LateXu (odkomentujte, chcete-li)
% %%% Uk�zka pou�it� n�kter�ch konstrukc� LaTeXu

\subsection{Uk�zka \LaTeX{}u}
\label{ssec:ukazka}

V~t�to kr�tk� ��sti uk�eme pou�it� n�kolika z�kladn�ch konstrukc� \LaTeX{}u,
kter� by se v�m mohly p�i psan� pr�ce hodit.

T�eba odr�ky:

\begin{itemize}
\item Logo Matfyzu vid�me na obr�zku~\ref{fig:mff}.
\item Tato subsekce m� ��slo~\ref{ssec:ukazka}.
\item Odkaz na literaturu~\cite{lamport94}.
\end{itemize}

Druhy poml�ek:
�erveno-�ern� (kr�tk�),
strana 16--22 (st�edn�),
$45-44$ (minus),
a~toto je --- jak se asi dalo �ekat --- vlo�en� v�ta ohrani�en� dlouh�mi poml�kami.
(V�imn�te si, �e jsme za \verb|a| napsali vlnovku m�sto mezery: to aby se
tam nemohl rozd�lit ��dek.)

\newtheorem{theorem}{V�ta}
\newtheorem*{define}{Definice}	% Definice ne��slujeme, proto "*"

\begin{define}
{\sl Strom} je souvisl� graf bez kru�nic.
\end{define}

\begin{theorem}
Tato v�ta neplat�.
\end{theorem}

\begin{proof}
Neplatn� v�ty nemaj� d�kaz.
\end{proof}

\begin{figure}
	\centering
	\includegraphics[width=30mm]{logo.eps}
	\caption{Logo MFF UK}
	\label{fig:mff}
\end{figure}


\chapter*{Z�v�r}
\addcontentsline{toc}{chapter}{Z�v�r}


%%% Seznam použité literatury
%%% Seznam použité literatury (bibliografie)
%%%
%%% Pro vytváření bibliografie používáme bibTeX. Ten zpracovává
%%% citace v textu (např. makro \cite{...}) a vyhledává k nim literaturu
%%% v souboru literatura.bib.
%%%
%%% Příkaz \bibliographystyle určuje, jakým stylem budou citovány odkazy
%%% v textu. V závorce je název zvoleného souboru .bst. Styly plainnat
%%% a unsrt jsou standardní součástí latexových distribucí. Styl czplainnat
%%% je dodáván s touto šablonou a bibTeX ho hledá v aktuálním adresáři.

\bibliographystyle{czplainnat}    %% Autor (rok) s českými spojkami
% \bibliographystyle{plainnat}    %% Autor (rok) s anglickými spojkami
% \bibliographystyle{unsrt}       %% [číslo]

\renewcommand{\bibname}{Seznam použité literatury}

%%% Vytvoření seznamu literatury. Pozor, pokud jste necitovali ani jednu
%%% položku, seznam se automaticky vynechá.

\bibliography{literatura}

%%% Kdybyste chtěli bibliografii vytvářet ručně (bez bibTeXu), lze to udělat
%%% následovně. V takovém případě se řiďte normou ISO 690 a zvyklostmi v oboru.

% \begin{thebibliography}{99}
%
% \bibitem{lamport94}
%   {\sc Lamport,} Leslie.
%   \emph{\LaTeX: A Document Preparation System}.
%   2. vydání.
%   Massachusetts: Addison Wesley, 1994.
%   ISBN 0-201-52983-1.
%
% \end{thebibliography}


%%% Tabulky v bakalářské práci, existují-li.
\chapwithtoc{Seznam tabulek}

%%% Použité zkratky v bakalářské práci, existují-li, včetně jejich vysvětlení.
\chapwithtoc{Seznam použitých zkratek}

%%% Přílohy k bakalářské práci, existují-li (různé dodatky jako výpisy programů,
%%% diagramy apod.). Každá příloha musí být alespoň jednou odkazována z vlastního
%%% textu práce. Přílohy se číslují.
\chapwithtoc{Přílohy}

\openright
\end{document}
